% GENERATED FILE. Edit canonical lessons, then run npm run generate:books.
% canonical-source-hash: fnv1a64:658c03aacd3d2e07
% canonical-lessons: FR-C416-vol, FR-C416-aerogare, FR-C416-enregistrement, FR-C416-soute, FR-C416-embarquement

\chapter{À l'aérogare}
\label{ch:fr-a-l-aerogare}
\hlchaptermodality{416}

\begin{chapteropening}
\textbf{By the end of this chapter:} \emph{I can check in for a flight: the flight, the airport building, check-in, the hold and boarding.}

Complete the last lesson of chapter 416: I can check in for a flight: the flight, the airport building, check-in, the hold and boarding.
\end{chapteropening}

\section[le vol]{le vol — the flight}

\label{lesson:FR-C416-vol}



\begin{quote}
Before the new one: say the French for the tram, then the French for the scooter.
\end{quote}

\subsection*{You'll want to know: le vol}
\textbf{le vol} — \textquotedblleft{}the flight\textquotedblright{}.

\begin{quote}
Notre vol part à sept heures du matin.
\end{quote}

\begin{quote}
\emph{Our flight leaves at seven in the morning.}
\end{quote}

\emph{Le vol} is also a theft; that sense comes from \emph{voler} 'to steal'.

\begin{cousinweb}
\emph{vol} is built on \emph{voler} 'to fly' (EtymDB), which is inherited from Latin \emph{volāre}.
\end{cousinweb}

\begin{grammarlens}[title={What you've built}]
One more word for telling, explaining and describing at B1.
\end{grammarlens}

\subsection*{Your turn}
\begin{itemize}
\raggedright
  \item \emph{Say it:} \emph{le vol}
  \item \emph{Say it:} \emph{le vol}, once more
  \item \emph{Say it:} a sentence of your own with \emph{le vol}
  \item \emph{Recall:} say the French for the tram, then the French for the scooter, then say \emph{le vol} again
\end{itemize}

\begin{tcolorbox}[breakable,colback=teal!4,colframe=teal!35!black,title={Before you move on}]
What does le vol mean? (\textquotedblleft{}The flight\textquotedblright{}.) Say it once more.
\end{tcolorbox}
\section[l'aérogare]{l'aérogare — the airport building (for passengers)}

\label{lesson:FR-C416-aerogare}



\begin{quote}
Before the new one: say the French for the scooter, then the French for the flight.
\end{quote}

\subsection*{You'll want to know: l'aérogare}
\textbf{l'aérogare} — \textquotedblleft{}the airport building (for passengers)\textquotedblright{}.

\begin{quote}
La navette va de la gare jusqu'à l'aérogare.
\end{quote}

\begin{quote}
\emph{The shuttle runs from the railway station to the airport building.}
\end{quote}

A feminine noun: \emph{une aérogare}.

\begin{cousinweb}
\emph{aérogare} is built from \emph{aéro-}, from Greek \emph{a\'{\={e}}r} 'air', and \emph{gare}, the building for trains: a \emph{gare} for planes. EtymDB lists the word without a source line.
\end{cousinweb}

\begin{grammarlens}[title={What you've built}]
One more word for telling, explaining and describing at B1.
\end{grammarlens}

\subsection*{Your turn}
\begin{itemize}
\raggedright
  \item \emph{Say it:} \emph{l'aérogare}
  \item \emph{Say it:} \emph{l'aérogare}, once more
  \item \emph{Say it:} a sentence of your own with \emph{l'aérogare}
  \item \emph{Recall:} say the French for the scooter, then the French for the flight, then say \emph{l'aérogare} again
\end{itemize}

\begin{tcolorbox}[breakable,colback=teal!4,colframe=teal!35!black,title={Before you move on}]
What does l'aérogare mean? (\textquotedblleft{}The airport building (for passengers)\textquotedblright{}.) Say it once more.
\end{tcolorbox}
\section[l'enregistrement]{l'enregistrement — the check-in; the recording}

\label{lesson:FR-C416-enregistrement}



\begin{quote}
Before the new one: say the French for the flight, then the French for the airport building.
\end{quote}

\subsection*{You'll want to know: l'enregistrement}
\textbf{l'enregistrement} — \textquotedblleft{}the check-in; the recording\textquotedblright{}.

\begin{quote}
L'enregistrement ferme une heure avant le départ.
\end{quote}

\begin{quote}
\emph{Check-in closes an hour before departure.}
\end{quote}

A masculine noun. Also a recording of music.

\begin{cousinweb}
\emph{enregistrement} is built on \emph{enregistrer} with \emph{-ment} (EtymDB). \emph{Enregistrer} comes from \emph{registre} 'register', from Medieval Latin \emph{registrum}.
\end{cousinweb}

\begin{grammarlens}[title={What you've built}]
One more word for telling, explaining and describing at B1.
\end{grammarlens}

\subsection*{Your turn}
\begin{itemize}
\raggedright
  \item \emph{Say it:} \emph{l'enregistrement}
  \item \emph{Say it:} \emph{l'enregistrement}, once more
  \item \emph{Say it:} a sentence of your own with \emph{l'enregistrement}
  \item \emph{Recall:} say the French for the flight, then the French for the airport building, then say \emph{l'enregistrement} again
\end{itemize}

\begin{tcolorbox}[breakable,colback=teal!4,colframe=teal!35!black,title={Before you move on}]
What does l'enregistrement mean? (\textquotedblleft{}The check-in; the recording\textquotedblright{}.) Say it once more.
\end{tcolorbox}
\section[la soute]{la soute — the hold (of a plane or ship)}

\label{lesson:FR-C416-soute}



\begin{quote}
Before the new one: say the French for the airport building, then the French for the check-in.
\end{quote}

\subsection*{You'll want to know: la soute}
\textbf{la soute} — \textquotedblleft{}the hold (of a plane or ship)\textquotedblright{}.

\begin{quote}
Sa grande valise voyage dans la soute.
\end{quote}

\begin{quote}
\emph{Her big suitcase travels in the hold.}
\end{quote}

\emph{Un bagage en soute} — a checked bag.

\begin{cousinweb}
\emph{soute} was borrowed from Provençal \emph{sota} (EtymDB), which goes back to Latin \emph{subtus} 'underneath': the space below deck.
\end{cousinweb}

\begin{grammarlens}[title={What you've built}]
One more word for telling, explaining and describing at B1.
\end{grammarlens}

\subsection*{Your turn}
\begin{itemize}
\raggedright
  \item \emph{Say it:} \emph{la soute}
  \item \emph{Say it:} \emph{la soute}, once more
  \item \emph{Say it:} a sentence of your own with \emph{la soute}
  \item \emph{Recall:} say the French for the airport building, then the French for the check-in, then say \emph{la soute} again
\end{itemize}

\begin{tcolorbox}[breakable,colback=teal!4,colframe=teal!35!black,title={Before you move on}]
What does la soute mean? (\textquotedblleft{}The hold (of a plane or ship)\textquotedblright{}.) Say it once more.
\end{tcolorbox}
\section[l'embarquement]{l'embarquement — the boarding}

\label{lesson:FR-C416-embarquement}



\begin{quote}
Before the new one: say the French for the check-in, then the French for the hold.
\end{quote}

\subsection*{You'll want to know: l'embarquement}
\textbf{l'embarquement} — \textquotedblleft{}the boarding\textquotedblright{}.

\begin{quote}
L'embarquement commence à la porte douze.
\end{quote}

\begin{quote}
\emph{Boarding begins at gate twelve.}
\end{quote}

A masculine noun. \emph{La carte d'embarquement} — the boarding pass.

\begin{cousinweb}
\emph{embarquement} is built on \emph{embarquer} 'to board' (EtymDB lists it), from \emph{en-} and \emph{barque} 'boat', which goes back to Latin \emph{barca} (Lewis \& Short).
\end{cousinweb}

\begin{grammarlens}[title={What you've built}]
One more word for telling, explaining and describing at B1.
\end{grammarlens}

\subsection*{Your turn}
\begin{itemize}
\raggedright
  \item \emph{Say it:} \emph{l'embarquement}
  \item \emph{Say it:} \emph{l'embarquement}, once more
  \item \emph{Say it:} a sentence of your own with \emph{l'embarquement}
  \item \emph{Recall:} say the French for the check-in, then the French for the hold, then say \emph{l'embarquement} again
\end{itemize}

\begin{tcolorbox}[breakable,colback=teal!4,colframe=teal!35!black,title={Before you move on}]
What does l'embarquement mean? (\textquotedblleft{}The boarding\textquotedblright{}.) Say it once more.
\end{tcolorbox}
